\chapter{Implementation}
\label{ch:implementation}

\section{Technology Stack}
\label{sec:stack}

\TODO{Brief and factual. Justify only the choices that were genuinely
decisions, and skip the ones that were defaults --- a bachelor thesis loses
credit for pretending every dependency was deliberated.}

\section{The Grounded Tutor}
\label{sec:tutor}

This is the content-consumption counterpart to \cref{sec:pipeline}: it is what
the derived objects exist for.

\thesisfigurepage{f12-tutor-sequence}
  {The grounded lecture tutor. Groundedness is enforced at four distinct points
   along the request path.}
  {fig:tutor-sequence}

\subsection{Retrieval Scoping}
\label{sec:retrieval}

\TODO{Emphasise that the lecture scope is applied in the SQL \code{WHERE}
clause rather than as a post-filter in application code, and explain why that
matters: with a Python post-filter over a global nearest-neighbour scan, a
lecture's relevant slides silently stop appearing once the candidate window
fills with other lectures' content. The migration comment explains this; cite
it.}

\TODO{Also cover the query-embedding cache, and the anchoring of the current
slide at a synthetic similarity of zero --- the code comment explains that this
prevents refusal logic from confusing "we chose this slide as an anchor" with
"this slide is highly relevant". A small detail that shows care.}

\subsection{Anti-Hallucination Mechanisms}
\label{sec:anti-hallucination}

\TODO{Four layers, in order: input sanitisation, scoped retrieval, prompt
constraints, and post-hoc citation validation that drops any cited slide index
absent from the retrieved set. Be precise about what each one can and cannot
guarantee. Post-hoc validation catches invented slide numbers; it does not
verify that the answer's content is actually supported by the cited slide.}

\subsection{Prompt Injection via Uploaded Content}
\label{sec:injection}

\TODO{Connect back to \cref{sec:stakeholders}: because uploads are not
professor-only, a crafted PDF can reach every other student's tutor context.
The mitigation is prompt-level --- the retrieved context is explicitly labelled
untrusted --- plus a sanitiser that escapes delimiters and rewrites injection
phrasings. Then state the limitation honestly: prompt-level defences against
prompt injection are mitigations, not sound guarantees. Saying so demonstrates
judgement; claiming otherwise invites an easy objection.}

\section{Assessment Generation}
\label{sec:assessment}

\TODO{How quiz questions are produced inline by the pipeline rather than by a
separate user action, and the quality constraints encoded in the prompt:
one concept-testing question per slide, four plausible same-length distractors
reflecting real misconceptions, a preference for application over recall, and an
explicit ban on "all of the above" style options. These are pedagogical
decisions expressed as prompt constraints, which is worth naming as such.}

\section{Spaced Repetition and Exam Mode}
\label{sec:srs-impl}

\TODO{The card factory derives review cards from existing questions with no
further LLM call --- a cost decision worth stating. Cover SM-2 behind its stable
interface, and mock-exam sampling: coverage-first, weakness-weighted, with a
server-generated seed so a student cannot replay a favourable one. That last
detail is a nice example of an integrity requirement met in the sampling design
rather than bolted on.}

\section{Human-in-the-Loop Review}
\label{sec:review}

\thesisfigurepage{f13-batch-review}
  {The professor batch-review flow. The review step is corrective, not gating:
   content is student-visible from the moment parsing finishes.}
  {fig:batch-review}

\TODO{This is a first-order finding for a thesis about generating educational
content, so do not bury it. There is no draft/published state --- the source
comment says so verbatim --- meaning AI-extracted content reaches students
before any human inspects it, and the completion marker is stored only in
browser local storage. Additionally, the review flag is never cleared on save,
so a corrected slide stays flagged permanently. Discuss what the pipeline would
need for review to become a genuine gate, and what that would cost in latency
for the professor.}

\section{Authorisation}
\label{sec:authorisation}

\thesisfigure{f14-trust-boundaries}
  {Trust boundaries. Authorisation is enforced in the database, not the API
   layer.}
  {fig:trust-boundaries}

\TODO{Explain row-level security as the enforcement point, then the four policy
patterns. The important structural point is that Postgres combines permissive
policies with a logical OR, so effective visibility is the union of every
applicable policy rather than the intersection --- and that no test composes a
real endpoint against real policies, so the union behaviour remains unguarded
at the integration level. A documented, citable limitation.}

\subsection{Iterative Hardening}
\label{sec:hardening}

\thesisfigure{f15-security-arc}
  {Five sequential migrations hardening privilege escalation, each closing a
   vulnerability the previous step created or left open.}
  {fig:security-arc}

\TODO{\Cref{fig:security-arc} is the strongest security material available:
each migration header documents its own attack, so the sources are primary. Walk
the five steps and be explicit that step three (opening professor self-signup)
\emph{introduced} the vulnerability step four closes --- an honest account of
iterative hardening is more instructive than a tidy one.}

\TODO{Then present the one fully reproduced exploit in detail, because it is a
teaching example: as an unauthenticated role, a security-definer function with a
caller-supplied identifier and no ownership check leaked a victim's social
graph. The root cause was a four-way conjunction including Postgres's implicit
grant of execute permission to PUBLIC on function creation. It was reproduced
against a real database and is covered by regression tests.}

\TODO{Also include the negative result, which is worth as much: seven other
functions are callable by the unauthenticated role but return no rows, because
they are security-\emph{invoker} and row-level security therefore evaluates as
the caller. Anon-callable is not anon-readable. Presenting the contrast shows
you understood the mechanism rather than pattern-matching on a keyword.}
